\section{Positive windows and the noncritical raw inverse}
\label{sec:v46-window-transport}
The fixed-band theorem controls interval probabilities at the original
integer labels. We use those positive probabilities to control a
noncritical density and its first derivative. The comparison moves the
two endpoint contacts, preserves every section decision, and changes
the roof by a prescribed amount. The number of collision words never
enters the bound. This supplies a local-scale estimate for a regular
part of the signed inverse, rather than a certificate for reconstructing
that part at a bandwidth chosen from an unknown derivative norm.

Throughout this section $m$ is the collision count, $n$ is the return
index, and $1\le n\le m$. Set
\[
 \Pi_{n,k,m,R}=\{K_{n,R}=k,\ N_{n,R}=m\}.
\]
All source restrictions are taken under the original probability
$\nu_R^*$, without subsequent normalization. On this event the roof
$T_{n,R}$ is the length of the same $m$-flight word. Endpoint coordinates
$u,v$ are contact arclengths. Local arclength coordinates differ by
translations or changes of orientation; the gradients and divergences
below use their flat metric.

\subsection{Protection away from a critical point}
Use $q,d_j,e_j,e'_j$ of Section~\ref{sec:v45-critical-collars}. Call an
arbitrary regular $m$-flight source state $\varepsilon$-protected if its
incidences, section distances and nonincident clearances satisfy the
same inequalities as there. Normality of the endpoints is not required.
Denote this set by $\mathsf P_{m,R}(\varepsilon)$, and fix a sufficiently
small constant $\varepsilon_0>0$. All statements below use
$0<\varepsilon<\varepsilon_0$.

\begin{lemma}[Endpoint continuation for a protected regular source]
\label{lem:v46-regular-endpoints}
There are constants $r,C>0$, independent of $m,R,\varepsilon$, such
that every source state in $\mathsf P_{m,R}(\varepsilon)$ has an endpoint
square of radius $r\varepsilon^3$ on which its physical word and every
section decision persist. If $F$ is the reduced action and $w\,du\,dv$
is the section source density on that square, then
\begin{align}
 \nabla F&=(-p_0,p_m),&
 0<\nabla^2F&\le C I_2,&
 |\nabla\log w|&\le C\varepsilon^{-3}.
                                      \label{eq:v46-regular-chart}
\end{align}
All margins remain at least half their original lower bounds. The
lower Hessian bound may be taken as $\varepsilon I_2/(2R_+)$. In a
square whose endpoint tangential velocities have modulus at most
$1/2$, it may instead be taken as a fixed positive multiple of $I_2$.
Here $R_+=47/100$.
\end{lemma}
\begin{proof}
The internal matrix \eqref{eq:v45-contact-matrix}, its row-sum bound,
and the cancellation in the endpoint--internal derivative do not use
normality of either endpoint. In fact their endpoint factors $c_0,c_m$
are at most one. Thus \eqref{eq:v45-endpoint-influence} and
\eqref{eq:v45-summed-influence} hold at every protected regular source.
The continuation argument in Lemma~\ref{lem:v45-endpoint-collar}
therefore preserves the same half-margins on the stated square. For
each fixed word, endpoint injectivity glues the local continuations;
there is no new multiplicity in these coordinates.

The first variation gives the gradient. The Schur-complement argument
of Lemma~\ref{lem:v44-endpoint-curvature} bounds the reduced Hessian
below by the two endpoint curvature terms. The upper bound is the
endpoint principal block of the full contact Hessian, since the
eliminated block is positive. Flight lengths are bounded below by
$3/50$ and $R\ge9/20$, so this upper bound is independent of the
internal incidences and of word length.

The cross-density formula \eqref{eq:v45-cross-density} is also valid
without normal endpoints. It retains the factor $(4\pi R c)^{-1}$,
where $c=\nu(Y_R^*)$. The internal potential terms in its logarithmic
derivative are bounded by
$C\varepsilon^{-3}\sum_j q^{d_j/4}$, as before. Now the two endpoint
cosines are only bounded below by $\varepsilon/2$. Their differentials
are at most $C\varepsilon^{-1}(|du|+|dv|)$: incidence is a Lipschitz
function of adjacent contact positions, and their endpoint responses
satisfy \eqref{eq:v45-summed-influence}. Their logarithmic differentials
therefore cost at most $C\varepsilon^{-2}(|du|+|dv|)$, still within the
claimed budget. The other length terms cost at most
$C\varepsilon^{-1}(|du|+|dv|)$. This proves the last estimate.
Empty internal matrices require only the two endpoint terms.
\end{proof}

Let $\Theta:[0,\infty)\to[0,1]$ be a fixed nondecreasing smooth
function which is zero on $[0,1]$, one on $[2,\infty)$, and has bounded
derivative. At contact $j$ let $s_j$ be the distance to
$\partial Y_R^*$ in the fixed collision cylinder, on either side of
that boundary. At flight $j$ let $\Delta_j$ be its least clearance
from a nonincident disk. Distances and clearances may be capped at one.
Only a fixed finite collection of disks can affect that capped
clearance. Define the source guard
\begin{equation}\label{eq:v46-source-guard}
 \Xi_{m,R}^{\varepsilon}=
 \prod_{j=0}^{m}\Theta\!\left(\frac{c_j}{\varepsilon e_j}\right)
                       \Theta\!\left(\frac{s_j}{\varepsilon e_j}\right)
 \prod_{j=0}^{m-1}\Theta\!\left(\frac{\Delta_j}{\varepsilon e'_j}\right).
\end{equation}
Set it to zero on nonregular states. This is a restriction of the
physical source, not a change of its collision record or return index.

\begin{lemma}[A guard with a word-length-independent derivative]
\label{lem:v46-guard-derivative}
The guard is between zero and one, is one on
$\mathsf P_{m,R}(2\varepsilon)$, and can be nonzero only on
$\mathsf P_{m,R}(\varepsilon)$. On regular endpoint charts it is
locally Lipschitz and
\begin{equation}\label{eq:v46-guard-gradient}
 |\nabla\Xi_{m,R}^{\varepsilon}|\le C\varepsilon^{-2}
                 \quad\hbox{almost everywhere}.
\end{equation}
For each fixed $m$ its zero extension vanishes near every boundary of
its physical word and of every section-decision domain. No constants
in \eqref{eq:v46-guard-gradient} grow with $m$.
\end{lemma}
\begin{proof}
The first assertions follow directly from \eqref{eq:v46-source-guard}.
Where a term in its derivative is nonzero, all the other nonzero
factors impose their $\varepsilon$-margins. Adjacent-contact formulas
and \eqref{eq:v45-endpoint-influence} consequently give
\[
 |\nabla c_j|+|\nabla s_j|
            \le C\varepsilon^{-1}q^{3d_j/4}.
\]
A corresponding bound holds for $\Delta_j$ with the flight distance
index; taking the minimum of finitely many Lipschitz clearances does
not increase their largest Lipschitz constant. Dividing a differentiated
factor by $\varepsilon e_j$ costs at most
$C\varepsilon^{-2}q^{d_j/2}$. At most two contacts occur at each
distance from the endpoints, and the geometric series is finite.
This proves \eqref{eq:v46-guard-gradient}. At points where a factor is
zero the same conclusion holds almost everywhere by the Lipschitz
product rule. For fixed $m$ each lower margin is strictly positive;
hence the zero extension has no boundary jump. Distances at rectangle
corners need not be smooth: their weak first derivatives suffice here.
\end{proof}

\subsection{A positive flow-box comparison}
For $h>0$ define the finite-time local-window budget
\begin{equation}\label{eq:v46-window-budget}
 \mathcal Q_m(h)=\frac{m^2}{2h}
   \sup_{R,n,k,t}\nu_R^*
       \bigl(\Pi_{n,k,m,R}\cap\{|T_{n,R}-t|<h\}\bigr).
\end{equation}
Theorem~\ref{thm:v41-uniform-transition} and the boundedness of its
transition kernel give
\begin{equation}\label{eq:v46-window-budget-limit}
 \limsup_{m\to\infty}\mathcal Q_m(h)\le C_0
                         \quad\hbox{for each fixed }h>0,
\end{equation}
with $C_0$ independent of $h$. The window remains fixed in this limit.
This upper bound uses all arithmetic branches and all targets; it uses
neither a positive microscopic denominator nor a trivial residue.

\begin{lemma}[Transport from a roof level to a positive window]
\label{lem:v46-positive-flow-box}
Fix $0<\delta<1/10$. On the protected source with
$|\nabla F|\ge\delta$, put $V=\nabla F/|\nabla F|^2$. For a fixed
sufficiently small $c_1>0$, set
\begin{equation}\label{eq:v46-flow-width}
 h_{\varepsilon,\delta}=c_1\varepsilon^3\delta^2.
\end{equation}
The flow of $V$ exists for $|s|\le h_{\varepsilon,\delta}$ from every
such point, changes $F$ to $F+s$, and keeps its original exact labels.
Its physical density Jacobian is between $e^{-1}$ and $e$.
Consequently, for every measurable subset $A$ of that protected
noncritical source, its roof density satisfies
\begin{equation}\label{eq:v46-slice-window}
 \rho_A(t)\le \frac{e}{2h_{\varepsilon,\delta}}
       \nu_R^*\bigl(\Pi_{n,k,m,R}\cap
                    \{|T_{n,R}-t|<h_{\varepsilon,\delta}\}\bigr)
\end{equation}
almost everywhere. The inequality also holds when $A$ is the union
of such subsets over all physical words with the same exact labels.
\end{lemma}
\begin{proof}
As long as $|\nabla F|\ge\delta/2$, the flow speed is at most
$2/\delta$. Its endpoint displacement is at most
$2h_{\varepsilon,\delta}/\delta$. Lemma~\ref{lem:v46-regular-endpoints}
and a first-exit argument show that this stays inside the endpoint
square and keeps the gradient above $\delta/2$, after reducing $c_1$.
Indeed $h_{\varepsilon,\delta}\le c_1\varepsilon^3\delta$ and
$h_{\varepsilon,\delta}\le c_1\delta^2$. All half-margins and all
section decisions are retained. The identity $VF=1$ gives the exact
roof translation.

Writing divergence relative to $w\,du\,dv$ gives
\[
 \operatorname{div}_w V
   =\operatorname{div}V+\nabla\log w\cdot V,
 \qquad
 |\operatorname{div}_w V|
       \le C(\delta^{-2}+\varepsilon^{-3}\delta^{-1}).
\]
This follows by differentiating $V$ and using the Hessian and
log-density bounds. Multiplication by
$h_{\varepsilon,\delta}$ gives at most
$Cc_1(\varepsilon^3+\delta)$. The Liouville formula for the flow
therefore bounds $w(\Phi_s z)|\det D\Phi_s(z)|/w(z)$ between
$e^{-1}$ and $e$.

On a regular level $F=t$, the flow map from that level times
$(-h_{\varepsilon,\delta},h_{\varepsilon,\delta})$ is injective.
Equality of two images first forces equality of the flow times, by
the roof value, and then equality of the starting points, by uniqueness
of the flow. Coarea and the density Jacobian give its source mass at
least
\[
 2h_{\varepsilon,\delta}e^{-1}
       \int_{A\cap\{F=t\}}\frac{w}{|\nabla F|}\,d\mathcal H^1.
\]
Its image lies in the positive event on the right of
\eqref{eq:v46-slice-window}. Different physical words remain disjoint
throughout the flow. Summing their mass inequalities thus proves the
same bound for their union, without a multiplicity or word-count
factor. Coarea proves the statement for almost every level.
\end{proof}

\subsection{A first-derivative budget for the actual regular density}
We specify the weighted class. An insertion $a$ is admissible with
bounds $(M,L)$ if $|a|\le M$ and, on each regular endpoint chart where
the $\varepsilon/2$-margins hold, it has weak endpoint gradient at most
$L$. Bounds on a complex gradient mean the sum of the bounds for its
real and imaginary parts. No derivative across a physical singularity
is requested. Products $a_0(x)a_1(T_R^m x)$ of uniformly bounded
Lipschitz collision-state functions are admissible with a common
$(M,L)$: the endpoint-state maps are $(u,-\partial_uF)$ and
$(v,\partial_vF)$, up to the fixed arclength conversion, and their
first derivatives are bounded by \eqref{eq:v46-regular-chart}.
This assertion does not include an arbitrary measurable path selector.

Put $H_\delta=\Theta(|\nabla F|/\delta)$ on each word, and define
\begin{equation}\label{eq:v46-regular-source}
 g^{\varepsilon,\delta,a}_{n,k,m,R}(t)\,dt
  =(T_{n,R})_*
       (\one_{\Pi_{n,k,m,R}}a\Xi_{m,R}^{\varepsilon}
                                      H_\delta\,\nu_R^*).
\end{equation}
Thus both the exact labels and the observed roof are unchanged.

\begin{theorem}[Noncritical density and derivative at the local scale]
\label{thm:v46-noncritical-density}
For every fixed $\varepsilon,\delta$ and every admissible insertion,
$g=g^{\varepsilon,\delta,a}_{n,k,m,R}$ belongs to $W^{1,\infty}(\R)$.
There is an absolute constant $C$ such that, for every $m\ge1$,
\begin{align}
 m^2\|g\|_\infty
   &\le CM\mathcal Q_m(h_{\varepsilon,\delta}),\notag\\
 m^2\|g'\|_\infty
   &\le C\mathcal T_{\varepsilon,\delta}(M,L)
                   \mathcal Q_m(h_{\varepsilon,\delta}),
                                      \label{eq:v46-regular-density-budget}\\
 \mathcal T_{\varepsilon,\delta}(M,L)
   &=L\delta^{-1}+M(\varepsilon^{-3}\delta^{-1}+\delta^{-2}).\notag
\end{align}
The bounds are uniform over $R,n,k$ and the whole stated insertion
class. In particular the two normalized norms have finite limsups,
with no factor exponential in the collision count.
\end{theorem}
\begin{proof}
Let $A$ be the support of
$a\Xi_{m,R}^{\varepsilon}H_\delta$ and of its weak first derivatives,
ignoring null sets. Every such source point is
$\varepsilon$-protected and has $|\nabla F|\ge\delta$.
The first assertion in \eqref{eq:v46-regular-density-budget} follows
from domination by $M\rho_A$ and
Lemma~\ref{lem:v46-positive-flow-box}.

For a compactly supported smooth roof test $\psi$, integration by
parts on the regular endpoint domains gives
\begin{align}
 \int g(t)\psi'(t)\,dt
  &=\sum_{\rm words}\int a\Xi^{\varepsilon}H_\delta w
                                  V\cdot\nabla(\psi\circ F)\,du\,dv\notag\\
  &=-\sum_{\rm words}\int
      \operatorname{div}(a\Xi^{\varepsilon}H_\delta wV)
                                            \psi(F)\,du\,dv.
                                      \label{eq:v46-density-derivative}
\end{align}
The guard vanishes near each physical or section-decision boundary,
and $H_\delta$ vanishes near every critical point. For fixed $m$ there
are finitely many word domains. Arclength seams cancel as coordinate
seams on the contact torus; they are not source boundaries.
Lipschitz approximation justifies the same formula for the stated
weak derivatives. There is no omitted boundary measure.

Almost everywhere on the nonzero integrand,
\[
 \bigl|\operatorname{div}(a\Xi^{\varepsilon}H_\delta wV)\bigr|
       \le C\mathcal T_{\varepsilon,\delta}(M,L)w\one_A.
\]
Indeed the insertion derivative costs $L/\delta$, the guard derivative
costs $CM\varepsilon^{-2}/\delta$, the gradient cutoff costs
$CM/\delta^2$, and the weighted divergence costs
$CM(\delta^{-2}+\varepsilon^{-3}\delta^{-1})$.
The second budget in \eqref{eq:v46-regular-density-budget} now follows
from \eqref{eq:v46-density-derivative} and the same positive slice
comparison. These identities identify the distributional derivative
as a bounded density, proving $W^{1,\infty}$, including the zero
extension outside the physical roof range. Finally use
\eqref{eq:v46-window-budget-limit} for the fixed flow width.
\end{proof}

Use exactly the real even kernel $K_B$ of
Section~\ref{sec:v43-full-source-inversion}. It is Schwartz,
$K_B(t)=BK_1(Bt)$, and $\int K_B=1$. Put
$\kappa_1=\int |s|\,|K_1(s)|\,ds<\infty$.

\begin{corollary}[A common-band estimate for the noncritical inverse]
\label{cor:v46-noncritical-unsmoothing}
For all $B>0$,
\begin{equation}\label{eq:v46-noncritical-unsmoothing}
 m^2\|g-K_B*g\|_\infty
 \le\frac{C\kappa_1}{B}\,
       \mathcal T_{\varepsilon,\delta}(M,L)
                         \mathcal Q_m(h_{\varepsilon,\delta}).
\end{equation}
Consequently this signed correction tends to zero at the $m^{-2}$
scale for every $B=B_m\to\infty$, with $\varepsilon,\delta$ fixed.
No upper bound on the rate of growth of $B_m$ is required.
\end{corollary}
\begin{proof}
The continuous representative of a $W^{1,\infty}$ function satisfies
$|g(t)-g(t-s)|\le\|g'\|_\infty|s|$. Integrate this inequality against
$|K_B(s)|$ and apply \eqref{eq:v46-regular-density-budget}.
The spectral input has been used only at the fixed positive window
$h_{\varepsilon,\delta}$. The bandwidth in this corollary belongs to
an approximate-identity inequality for a proved density derivative;
it is not substituted into a fixed-band spectral theorem.
\end{proof}
